\documentclass[10pt]{article}
\usepackage[T1]{fontenc}
\usepackage[utf8]{inputenc}
\usepackage[french]{babel}
\usepackage{amsmath,amssymb,mathrsfs,mathpazo}
\usepackage{geometry}
\geometry{a4paper,margin=2cm}
\usepackage{graphicx,array,tabularx,tabularray,colortbl,multirow,multicol,booktabs,enumitem,pifont,ulem,fancybox,diagbox,cancel}
\usepackage{pstricks,pstricks-add,pst-plot,pst-tree,pst-node,pst-text,pst-func,pst-math,pst-bspline,pst-3dplot,pst-eucl,pst-solides3d}
\let\pscalculate\undefined
\usepackage{pst-bezier}
\usepackage{tikz,pgfplots,tkz-tab}
\usetikzlibrary{arrows,arrows.meta,patterns,calc,positioning,trees,decorations.markings,shapes.geometric,shapes.misc,shapes.arrows,intersections,decorations.pathmorphing,decorations.pathreplacing,backgrounds,shadows,babel}
\pgfplotsset{compat=1.18}
\usepgfplotslibrary{fillbetween}
\usepackage[np]{numprint}
\usepackage{listings,makecell,pas-tableur,twoopt}
\let\barre\undefined
\usepackage[pstricks]{tablvar}
\definecolor{darkgreen}{RGB}{0,100,0}
\providecommand{\red}{\color{red}}
\providecommand{\blue}{\color{blue}}
\providecommand{\green}{\color{green}}
\providecommand{\black}{\color{black}}
\providecommand{\white}{\color{white}}
\providecommand{\gray}{\color{gray}}
\providecommand{\cyan}{\color{cyan}}
\providecommand{\magenta}{\color{magenta}}
\providecommand{\yellow}{\color{yellow}}
\providecommand{\violet}{\color{violet}}
\providecommand{\orange}{\color{orange}}
\providecommand{\euro}{\texteuro}
\providecommand{\EUR}{\texteuro}
\providecommand{\ieme}{\textsuperscript{e}}
\providecommand{\degres}{\ensuremath{^\circ}}
\providecommand{\degree}{\ensuremath{^\circ}}
\providecommand{\vv}[1]{\overrightarrow{#1}}
\providecommand{\Coord}[1]{#1}
\providecommand{\decosix}{}
\providecommand{\starredbullet}{\ensuremath{\bullet}}
\pagestyle{empty}
\setlength{\parindent}{0pt}
\usepackage{adjustbox}\DeclareUnicodeCharacter{2236}{:}\DeclareUnicodeCharacter{2190}{\ensuremath{\leftarrow}}
\providecommand{\R}{\mathbb{R}}
\providecommand{\N}{\mathbb{N}}
\providecommand{\Z}{\mathbb{Z}}
\providecommand{\Q}{\mathbb{Q}}
\providecommand{\D}{\mathbb{D}}
\providecommand{\C}{\mathbb{C}}
\providecommand{\e}{{\mathrm{e}}}
\providecommand{\vect}[1]{\overrightarrow{#1}}
\providecommand{\vectt}[1]{\overrightarrow{\mathrm{#1}}}
\providecommand{\vvt}[1]{\overrightarrow{\mathrm{#1}}}
\providecommand{\np}[2][]{{#2~#1}}
\providecommand{\vv}[1]{{\overrightarrow{#1}}}
\renewcommand{\d}{{\mathrm{d}}}
\providecommand{\degres}{{^\circ}}
\providecommand{\degre}{{^\circ}}
\providecommand{\blue}{\color{blue}}
\providecommand{\red}{\color{red}}
\providecommand{\enskip}{\quad}
\providecommand{\llbracket}{\lbrack\!\lbrack}
\providecommand{\rrbracket}{\rbrack\!\rbrack}
\providecommand{\euro}{{\text{€}}}
\providecommand{\no}{{\text{n°}}}
\providecommand{\interval}[2]{{[#1;#2]}}
\providecommand{\intervFO}[2]{{[#1;#2[}}
\providecommand{\dotfill}{\ldots}

\providecommand{\ds}{\displaystyle}
\providecommand{\textding}[1]{#1}
\providecommand{\Oij}{\(\mathrm{O};\vec{\imath},\vec{\jmath}\)}
\providecommand{\Oijk}{\(\mathrm{O};\vec{\imath},\vec{\jmath},\vec{k}\)}
\providecommand{\Ouv}{\(\mathrm{O};\vec{u},\vec{v}\)}
\providecommand{\intervFF}[2]{\ensuremath{[#1;#2]}}
\renewcommand{\euro}{\ifmmode\text{\texteuro}\else\texteuro\fi}
\DeclareUnicodeCharacter{00B0}{\ifmmode^{\circ}\else\textdegree\fi}
\DeclareUnicodeCharacter{20AC}{\ifmmode\text{\texteuro}\else\texteuro\fi}
% Énoncés : Éducation nationale. Transcription et corrigés : APMEP, Denis Vergès et contributeurs. Découpage et classement : CoopMaths, Sébastien Lozano, Éric Elter et contributeurs. Adaptation de la présentation : Maths Explorer.
% CC BY-SA : https://coopmaths.fr/www/about/
% https://coopmaths.fr/alea/static/bac/2024/tex/bac_2024_09_sujet1_polynesie_4.tex
\begin{document}
\noindent\textit{APMEP / CoopMaths. CC BY-SA. Présentation adaptée par Maths Explorer.}\par\medskip
\medskip

\begin{minipage}{0.5\linewidth}
On considère un cube ABCDEFGH et l'espace est rapporté au repère orthonormal
$\left(\text{A}~;~\vect{\text{AB}},~\vect{\text{AD}},~\vect{\text{AE}}\right)$.

Pour tout réel $m$ appartenant à l'intervalle [0~;~ 1], on considère les points $K$ et $L$ de coordonnées:

\begin{center}$K(m~;~0~;~0)$\quad et \quad $L(1 - m~;~1~;~1)$.\end{center}
\end{minipage}\hfill
\begin{minipage}{0.45\linewidth}
\psset{unit=0.85cm}
\begin{pspicture}(6.6,6.6)
\pspolygon(0.4,1)(4.2,0.5)(4.2,4.8)(0.4,5.3)
\psline(4.2,0.5)(6.6,1.4)(6.6,5.7)(4.2,4.8)
\psline(6.6,5.7)(2.8,6.2)(0.4,5.3)
\psline(5,5.9)(0.4,5.3)(1.6,0.82)
\psline[linestyle=dashed](5,5.9)(6.6,1.4)(1.6,0.84)
\psline[linestyle=dashed](0.4,1)(2.8,1.9)(6.6,1.4)
\psline[linestyle=dashed](2.8,1.9)(2.8,6.2)
\uput[dl](0.4,1){A} \uput[d](4.2,0.5){B} \uput[dr](6.6,1.4){C} \uput[ul](2.8,1.9){D}
\uput[ul](0.4,5.3){E} \uput[dr](4.2,4.8){F} \uput[ur](6.6,5.7){G} \uput[u](2.8,6.2){H}
\uput[dl](1.6,0.84){K} \uput[u](5,5.9){L}
\end{pspicture}
\end{minipage}

\medskip

\begin{enumerate}
\item Donner les coordonnées des points E et C dans ce repère.
\item Dans cette question, $m = 0$. Ainsi, le point L(1~;~1~;~1) est confondu avec le point G, le point K(0~;~0~;~0) est confondu avec le point A et le plan (LEK) est donc le plan (GEA).
	\begin{enumerate}
		\item Justifier que le vecteur $\vect{\text{DB}}\begin{pmatrix}1\\-1\\0\end{pmatrix}$ est normal au plan (GEA).
		\item Déterminer une équation cartésienne du plan (GEA).
	\end{enumerate}
\end{enumerate}

On s'intéresse désormais à la nature de C$K$E$L$ en fonction du paramètre $m$.

\begin{enumerate}[resume]
\item Dans cette question, $m$ est un réel quelconque de l'intervalle [0~;~1].
	\begin{enumerate}
		\item Démontrer que C$K$E$L$ est un parallélogramme.
		\item Justifier que $\vect{K\text{C}} \cdot \vect{K\text{E}} = m(m - 1)$.
		\item Démontrer que C$K$E$L$ est un rectangle si, et seulement si, $m = 0$ ou $m = 1$.
	\end{enumerate}
\item Dans cette question, $m = \dfrac12$.Ainsi, L a pour coordonnées 
$\left(\dfrac12~;~1~;~1\right)$ et K a pour coordonnées $\left(\dfrac12~;~0~;~0\right)$.
	\begin{enumerate}
		\item Démontrer que le parallélogramme CKEL est alors un losange.
		\item À l'aide de la question \textbf{3.~b.}, déterminer une valeur approchée au degré près de la mesure de l'angle $\widehat{\text{CKE}}$.
	\end{enumerate}
\end{enumerate}

\end{document}
